\documentclass[12pt]{article}
\usepackage[utf8]{inputenc}
\usepackage[T1]{fontenc}
\usepackage[margin=1in]{geometry}
\usepackage{setspace}
\usepackage{xcolor}
\usepackage[colorlinks=true, urlcolor=blue!60!black, linkcolor=blue!60!black]{hyperref}
\setlength{\parindent}{2em}
\setlength{\parskip}{1.25em}
\renewcommand{\baselinestretch}{1.0}
\begin{document}
\begin{center}
{\LARGE \textit{Church of the Lukumi Babalu Aye, Inc. v. City of Hialeah}, 508 U.S. 520 (1993)\par}
\vspace{0.8em}
{\large Case Brief\par}
\end{center}
\vspace{1em}

\section{Memory Jogger}

Municipal ordinances gerrymandered to ban Santeria ritual animal sacrifice while exempting almost all secular and kosher animal killings are neither neutral nor generally applicable under \textit{Employment Division v. Smith}, and therefore trigger---and fail---strict scrutiny under the Free Exercise Clause.

\section{Detailed Case Facts}

Petitioner Church of the Lukumi Babalu Aye, Inc. (a Florida not-for-profit corporation organized in 1973) and its president and priest, Ernesto Pichardo (holding the title of \textit{Italero}, the second highest rank in the faith), practiced Santeria (``the way of the saints''), a nineteenth-century Afro-Caribbean syncretic religion formed when enslaved Yoruba people brought from West Africa to Cuba fused their traditional worship of spirits (\textit{orishas}) with Roman Catholic iconography (\href{https://www.courtlistener.com/opinion/112881/church-of-the-lukumi-babalu-aye-inc-v-city-of-hialeah/}{CourtListener}; \href{https://supreme.justia.com/cases/federal/us/508/520/}{Justia}). Persecuted in Cuba and practiced largely in secret, Santeria came to South Florida primarily with Cuban exiles, numbering at least 50,000 adherents in the region by the late 1980s.

A central element of Santeria devotion is animal sacrifice, through which adherents nurture personal relationships with the \textit{orishas} during birth, marriage, death, healing, initiation, and annual celebration rites. Sacrificed animals---including chickens, pigeons, doves, ducks, guinea pigs, goats, sheep, and turtles---are killed by cutting the carotid arteries in the neck, and the meat is cooked and eaten after all rituals except healing and death rites. In April 1987, the Church leased land in the City of Hialeah, Florida, and announced plans to establish a public house of worship, school, cultural center, and museum, intending to bring Santeria and its sacrificial rituals into the open. By early August 1987, the Church had obtained all required zoning, inspection, and licensing approvals.

The prospect of a Santeria church triggered intense public opposition in Hialeah. On June 9, 1987, the Hialeah City Council convened an emergency public session marked by overt hostility from council members, city officials, and a cheering audience that taunted Pichardo (\href{https://www.law.cornell.edu/supct/html/91-948.ZO.html}{Cornell LII}). Council members declared that Santeria adherents ``are in violation of everything this country stands for,'' asked ``[w]hat can we do to prevent the Church from opening?'', and questioned why a religion outlawed in Cuba should be brought to the United States; the Hialeah Police Department chaplain exhorted the council ``not to permit this Church to exist.'' Over the next three months, the council unanimously passed two resolutions and four substantive ordinances (each ordinance punishable by up to a \$500 fine, 60 days in jail, or both):

\begin{itemize}
  \item \textbf{Resolution 87-66 (June 9, 1987):} Noted residents' ``concern'' that ``certain religions may propose to engage in practices which are inconsistent with public morals, peace or safety'' and reiterated the City's commitment to prohibiting all such acts.
  \item \textbf{Ordinance 87-40 (June 9, 1987):} Incorporated Florida's animal-cruelty statute (Fla. Stat. ch. 828, including \S~828.12 punishing ``[w]hoever \dots{} unnecessarily or cruelly \dots{} kills any animal''), which the Florida Attorney General subsequently advised (Fla. Op. Atty. Gen. 87-56) prohibited religious animal sacrifice as ``unnecessary.''
  \item \textbf{Resolution 87-90 (August 11, 1987):} Declared city policy ``to oppose the ritual sacrifices of animals'' and warned that practitioners ``will be prosecuted.''
  \item \textbf{Ordinance 87-52 (September 8, 1987):} Defined ``sacrifice'' as ``to unnecessarily kill, torment, torture, or mutilate an animal in a public or private ritual or ceremony not for the primary purpose of food consumption,'' and prohibited possessing, sacrificing, or slaughtering animals for any ritual intending to use them for food, while exempting licensed establishments slaughtering animals specifically raised for food (thereby exempting kosher slaughter).
  \item \textbf{Ordinance 87-71 (September 22, 1987):} Directly outlawed sacrificing any animal within city limits, using the same definition of ``sacrifice'' as Ordinance 87-52.
  \item \textbf{Ordinance 87-72 (September 22, 1987):} Defined ``slaughter'' as killing animals for food and banned slaughter outside areas zoned for slaughterhouses, but exempted the slaughter or processing for sale of ``small numbers of hogs and/or cattle per week'' permitted under state law (Fla. Stat. \S~828.24(3)).
\end{itemize}

\section{Procedural History}

Following enactment of the ordinances, the Church and Pichardo filed suit under 42 U.S.C. \S~1983 in the United States District Court for the Southern District of Florida against the City of Hialeah and its mayor and city council members in their individual capacities, seeking declaratory, injunctive, and monetary relief for violations of, \textit{inter alia}, the Free Exercise Clause of the First Amendment (applicable to the States via the Fourteenth Amendment). The District Court granted summary judgment to the individual defendants on absolute legislative immunity (688 F. Supp. 1522 (S.D. Fla. 1988)).

After a nine-day bench trial on the remaining claims, the District Court entered judgment for the City (723 F. Supp. 1467 (S.D. Fla. 1989)). Although the court acknowledged that ``the ordinances are not religiously neutral'' and were ``prompted'' by the Church's arrival, it concluded that their purpose was to end animal sacrifice generally, and that four compelling governmental interests---preventing public-health risks, protecting children from emotional trauma, preventing animal cruelty, and enforcing slaughterhouse zoning---outweighed the religious burden because Santeria's secrecy made narrower regulations unenforceable (\href{https://supreme.justia.com/cases/federal/us/508/520/}{Justia}).

The United States Court of Appeals for the Eleventh Circuit affirmed in a one-paragraph \textit{per curiam} opinion (936 F.2d 586 (11th Cir. 1991)), declining to rely on the child-welfare interest and declining to evaluate the intervening decision in \textit{Employment Division v. Smith}, 494 U.S. 872 (1990), on the ground that the District Court had already applied an arguably stricter balancing standard.

The Supreme Court granted certiorari (503 U.S. 935 (1992)), heard oral argument on November 4, 1992, and unanimously reversed on June 11, 1993 (\href{https://www.courtlistener.com/opinion/112881/church-of-the-lukumi-babalu-aye-inc-v-city-of-hialeah/}{CourtListener}).

\section{Judicial Votes}

The Court voted unanimously (9--0) to reverse and invalidate all four ordinances, though the Justices divided over legislative-motive methodology and the continued validity of \textit{Employment Division v. Smith} (\href{https://caselaw.findlaw.com/court/us-supreme-court/508/520.html}{FindLaw}):

\begin{itemize}
  \item \textbf{Majority / Opinion of the Court (Justice Kennedy):}
  \begin{itemize}
    \item \textbf{Parts I, III, and IV:} Joined by Chief Justice Rehnquist and Justices White, Stevens, Scalia, Souter, and Thomas (7 Justices).
    \item \textbf{Part II-B (General Applicability):} Joined by Chief Justice Rehnquist and Justices White, Stevens, Scalia, and Thomas (6 Justices).
    \item \textbf{Parts II-A-1 and II-A-3 (Neutrality---Text, Operation, and Summary):} Joined by Chief Justice Rehnquist and Justices Stevens, Scalia, and Thomas (5 Justices).
    \item \textbf{Part II-A-2 (Legislative History / Motive Analysis):} Joined only by Justice Stevens (2 Justices; plurality/opinion of Kennedy, J.).
  \end{itemize}
  \item \textbf{Concurrences in Part and in the Judgment (2 separate opinions):}
  \begin{itemize}
    \item Justice Scalia filed an opinion concurring in part and concurring in the judgment, joined by Chief Justice Rehnquist (joining all of Kennedy's opinion except Part II-A-2).
    \item Justice Souter filed an opinion concurring in part and concurring in the judgment (joining only Parts I, III, and IV).
  \end{itemize}
  \item \textbf{Concurrence in the Judgment (1 opinion):}
  \begin{itemize}
    \item Justice Blackmun filed an opinion concurring in the judgment, joined by Justice O'Connor (joining none of Kennedy's opinion).
  \end{itemize}
  \item \textbf{Dissents:} None.
\end{itemize}

\section{Holding}

\textbf{Judgment:} Reversed; Hialeah's four challenged ordinances are unconstitutional and void.

\textbf{Rule Announced:} Under the Free Exercise Clause and \textit{Employment Division v. Smith}, a law that burdens religious practice need not satisfy strict scrutiny if it is neutral and of general applicability. Conversely, if a law is either not neutral (having as its object the infringement or restriction of practices because of their religious motivation, whether overtly or through a ``religious gerrymander'') or not generally applicable (selectively pursuing governmental interests only against religiously motivated conduct while leaving analogous secular conduct unprohibited), it must undergo ``the most rigorous of scrutiny'': it must advance interests ``of the highest order'' and be narrowly tailored in pursuit of those interests. Laws that are substantially underinclusive with respect to their asserted secular goals fail both general applicability and strict scrutiny.

\textbf{Holding vs. Dicta:} The Court's holding establishes that covertly gerrymandered, substantially underinclusive ordinances targeting a religious practice violate the Free Exercise Clause and fail strict scrutiny. As Justice Souter noted, the majority's restatement of the \textit{Smith} safe harbor for genuinely neutral, generally applicable laws was unnecessary to the judgment (and thus arguably dicta), because Hialeah's ordinances failed even the minimum nonpersecution baseline. Likewise, Justice Kennedy's Part II-A-2 endorsement of \textit{Arlington Heights} legislative-motive inquiry commanded only two votes (Kennedy and Stevens) and did not constitute a holding of the Court, and the Court expressly declined to define with precision the outer boundary of ``general applicability'' because Hialeah's laws fell ``well below the minimum standard.''

\section{Analysis of Opinions}

\subsection{Opinion of the Court (Justice Kennedy)}

Justice Kennedy began by confirming that Santeria is a protected religion and that its adherents' sincerity was undisputed, observing that religious beliefs need not be ``acceptable, logical, consistent, or comprehensible to others'' to merit First Amendment protection (\href{https://www.law.cornell.edu/supct/html/91-948.ZO.html}{Cornell LII}). Framing the inquiry under \textit{Smith}, Kennedy explained that neutrality and general applicability are ``interrelated,'' and that failing either triggers strict scrutiny.

In Part II-A-1 (for a 5-Justice majority), Kennedy held that facial neutrality is merely the minimum requirement and is not determinative: the Free Exercise Clause forbids ``subtle departures from neutrality'' and ``religious gerrymanders.'' Although the words ``sacrifice'' and ``ritual'' admit of secular meanings, their use alongside Resolution 87-66 and the interlocking definitions and exemptions of Ordinances 87-40, 87-52, and 87-71 proved that Santeria sacrifice was virtually the only conduct prohibited. By exempting killings for the ``primary purpose of food consumption'' and licensed food establishments (protecting kosher slaughter), and by interpreting Florida's ``unnecessary'' killing ban in Ordinance 87-40 to allow hunting, fishing, pest extermination, euthanasia, and live-rabbit greyhound training while deeming religious sacrifice ``unnecessary,'' Hialeah created a system of individualized governmental assessments that devalued religious reasons relative to secular ones. Moreover, the ordinances visited ``gratuitous restrictions'' far broader than needed for public health or humane slaughter (which could be addressed by organic-garbage disposal rules or humane-killing standards rather than banning Santeria sacrifices even inside licensed slaughterhouses). Ordinance 87-72, though broader on its face, functioned in tandem with the other enactments and had to fall with them.

In Part II-A-2 (joined only by Justice Stevens), Kennedy drew on equal protection doctrine (\textit{Arlington Heights v. Metropolitan Housing Development Corp.}, 429 U.S. 252 (1977)) to examine historical background, sequence of events, and contemporaneous statements by council members and city officials, concluding that the ordinances were enacted ``because of,'' not merely ``in spite of,'' their suppression of Santeria.

In Part II-B (for a 6-Justice majority), Kennedy held that all four ordinances failed general applicability because they were substantially underinclusive with respect to both asserted interests. On animal cruelty, Florida and Hialeah permitted fishing, hunting, pest extermination, and euthanasia; on public health, Hialeah left unregulated carcass disposal by hunters, garbage disposal by restaurants, and consumption of uninspected meat by hunters, fishermen, and home livestock owners. Ordinance 87-72 was facially underinclusive because it exempted commercial slaughter of ``small numbers of hogs and/or cattle.'' Thus, the ordinances imposed on Santeria worshippers a burden society was unwilling to impose upon itself.

In Parts III and IV (for a 7-Justice majority), Kennedy applied strict scrutiny and held that the ordinances failed both prongs: they were neither narrowly tailored (being simultaneously overbroad and underinclusive) nor backed by compelling interests, because ``a law cannot be regarded as protecting an interest `of the highest order' \dots{} when it leaves appreciable damage to that supposedly vital interest unprohibited.''

\subsection{Concurrence in Part and in the Judgment (Justice Scalia, joined by Chief Justice Rehnquist)}

Justice Scalia joined all of the Court's opinion except Part II-A-2 (\href{https://supreme.justia.com/cases/federal/us/508/520/}{Justia}). He first noted that ``neutrality'' and ``general applicability'' substantially overlap: neutrality primarily concerns facial religious classifications, while general applicability concerns facially neutral laws whose design, construction, or enforcement targets religious practice. Scalia refused to join Part II-A-2 because he rejected judicial inquiry into the subjective motives of lawmakers. Determining the singular motive of a multimember legislative body is ``virtually impossible,'' and the First Amendment's text prohibits laws that \textit{in effect} prohibit free exercise, not laws passed with impure thoughts. Even if every Hialeah council member had acted solely out of pure-hearted concern for animal welfare, the gerrymandered, underinclusive ordinances would still be invalid; conversely, an ineptly drafted law that failed to burden Santeria would not violate the Clause regardless of evil motive.

\subsection{Concurrence in Part and in the Judgment (Justice Souter)}

Justice Souter joined Parts I, III, and IV, agreeing that Hialeah's laws violated the noncontroversial ``nonpersecution principle,'' but declined to join Part II because it repeated \textit{Smith}'s rule in dicta (\href{https://supreme.justia.com/cases/federal/us/508/520/}{Justia}). Distinguishing ``formal neutrality'' (barring deliberate targeting) from ``substantive neutrality'' (requiring exemptions from formally neutral laws that substantially burden religion), Souter mounted a comprehensive critique of \textit{Smith}: he argued that \textit{Smith} rested on unpersuasive distinctions of \textit{Wisconsin v. Yoder} and the unemployment-compensation cases, was announced \textit{sua sponte} without full briefing, created an ``intolerable tension'' in precedent, and conflicted with both the text of the Free Exercise Clause and historical scholarship (citing Michael W. McConnell) on its original understanding. Souter urged the Court to reexamine \textit{Smith} in a future case where its validity was squarely presented.

\subsection{Concurrence in the Judgment (Justice Blackmun, joined by Justice O'Connor)}

Justice Blackmun, joined by Justice O'Connor, concurred only in the judgment (\href{https://supreme.justia.com/cases/federal/us/508/520/}{Justia}). Reaffirming his dissent in \textit{Smith}, Blackmun argued that \textit{Smith} wrongly reduced the Free Exercise Clause to a mere antidiscrimination rule rather than an affirmative substantive liberty. Under \textit{Sherbert v. Verner} and \textit{Yoder}, any law that substantially burdens religiously motivated conduct---intentionally or unintentionally---must be the least restrictive means of achieving a compelling state interest. Because a law that directly targets religion is ``by definition'' neither narrowly tailored nor backed by a genuine compelling interest, it ``\textit{ipso facto}'' fails strict scrutiny, making \textit{Lukumi} an easy case. Blackmun closed by emphasizing that a genuinely neutral, generally applicable animal-cruelty law would present a much harder question, noting the seriousness of animal-welfare interests raised by amici.

\section{Examples: Future Applications}

\begin{enumerate}
  \item \textbf{Same Side (Unconstitutional under \textit{Lukumi}):} A city council, reacting to a newly proposed mosque's plans to perform traditional \textit{halal} slaughter on-site during Eid al-Adha, enacts an ordinance banning ``ceremonial animal killing'' within city limits while exempting commercial abattoirs, recreational hunting, and backyard poultry slaughter for family dinners. Because the ordinance gerrymanders exemptions to target a minority religious practice and is grossly underinclusive regarding animal welfare and sanitation, it fails both neutrality and general applicability under \textit{Lukumi} and cannot survive strict scrutiny.

  \item \textbf{Same Side (Unconstitutional under \textit{Lukumi}):} During a public-health emergency, a state governor limits in-home religious gatherings and worship services to three households while allowing secular commercial venues---such as hair salons, retail stores, and restaurants---that pose comparable transmission risks to host larger crowds. Applying \textit{Lukumi}'s general-applicability principle (as the Supreme Court later did in \textit{Tandon v. Newsom}, 593 U.S. 61 (2021), and as analyzed in \href{https://lawreview.gmu.edu/print__issues/religious-and-secular-comparators/}{George Mason Law Review}), treating comparable secular activities more favorably than religious exercise strips the regulation of general applicability and triggers strict scrutiny.

  \item \textbf{Opposite Side (Constitutional under \textit{Lukumi} and \textit{Smith}):} A state enforces a longstanding, across-the-board criminal statute banning the possession or sale of endangered migratory birds without any secular exemptions for hunting, commercial trade, or scientific collection. A religious group that uses feathers from that endangered species in initiation rites seeks a constitutional exemption. Because the statute is facially and operationally neutral, contains no individualized-exemption mechanism, and is not underinclusive, it satisfies \textit{Smith} and \textit{Lukumi} and is subject only to rational-basis review under the Free Exercise Clause (absent statutory protection such as RFRA).

  \item \textbf{Opposite Side (Constitutional under \textit{Lukumi} and \textit{Smith}):} A city enacts a uniform sanitation ordinance requiring all persons and businesses---including restaurants, hunters, butcher shops, pet clinics, and houses of worship---to seal organic animal waste in leak-proof containers and place it in designated municipal waste receptacles rather than open alleys or parks. Even though the rule incidentally affects how Santeria adherents dispose of carcasses after healing and death rites, it is neutral, generally applicable, and directly targets the secular collateral harm of improper waste disposal---precisely the type of narrower regulation Justice Kennedy indicated Hialeah could have validly adopted.

  \item \textbf{Fence-Sitter (Hard Case):} A state civil-rights agency enforces a facially neutral public-accommodations law prohibiting sexual-orientation discrimination against a religious baker who declines to design a custom same-sex wedding cake; during administrative hearings, two commissioners make disparaging remarks comparing religious objections to slavery and the Holocaust, while the agency previously allowed three secular bakers to refuse cakes with anti-gay religious messages. As in \href{https://supreme.justia.com/cases/federal/us/584/16-111/}{\textit{Masterpiece Cakeshop, Ltd. v. Colorado Civil Rights Commission}}, 584 U.S. 617 (2018), the case sits on the knife-edge of \textit{Lukumi}'s unresolved fracture: under Justice Kennedy's Part II-A-2 approach, adjudicative hostility alone taints neutrality, whereas under Justice Scalia's concurrence, courts should look only to the law's objective text and uneven enforcement rather than officials' subjective statements.
\end{enumerate}

\section{Critique}

Although \textit{Lukumi} commanded a unanimous 9--0 judgment, it exposed deep theoretical fault lines in Free Exercise jurisprudence that continue to shape Supreme Court litigation (\href{https://fedsoc.org/case/church-of-the-lukumi-babalu-aye-inc-v-city-of-hialeah}{Federalist Society}; \href{https://fedsoc.org/case/fulton-v-city-of-philadelphia}{Federalist Society on \textit{Fulton}}).

\textbf{Logical and Doctrinal Weaknesses:} First, because Hialeah's ordinances were such an egregious, almost caricatured example of religious targeting, the Court expressly declined to define how much secular underinclusiveness destroys ``general applicability.'' That omission left lower courts divided for nearly three decades over whether a single secular exemption (the ``most-favored-nation'' theory) or only pervasive targeting triggers strict scrutiny---a doctrinal ambiguity that later drove contentious rulings in \textit{Fulton v. City of Philadelphia} (2021) and \textit{Tandon v. Newsom} (2021) (\href{https://lawreview.gmu.edu/print__issues/religious-and-secular-comparators/}{George Mason Law Review}). Second, Justice Kennedy's Part II-A-2 reliance on \textit{Arlington Heights} legislative-motive evidence received only two votes, creating lasting uncertainty about when hostile remarks by individual officials invalidate otherwise neutral laws or adjudications (an issue revived in \href{https://supreme.justia.com/cases/federal/us/584/16-111/}{\textit{Masterpiece Cakeshop}}). Third, as Justice Blackmun observed, applying full strict scrutiny after already finding purposeful religious persecution is arguably redundant: a law whose object is religious suppression can rarely, if ever, be justified by a compelling secular interest.

\textbf{Originalist and Textualist Perspectives:} Originalist and textualist scholars split along the very fault line between Justices Scalia and Souter. Textualist defenders of Scalia's concurrence praise his refusal to psychoanalyze multimember legislatures, arguing that the First Amendment's text (``Congress shall make no law \dots{} prohibiting the free exercise'') focuses on what a law objectively \textit{does} rather than why individual lawmakers voted for it. Conversely, originalist critics of \textit{Smith} (most prominently Michael W. McConnell, whose 1990 \textit{Harvard Law Review} article Justice Souter cited at length) contend that \textit{Lukumi}'s preservation of \textit{Smith}'s formal-neutrality framework falls short of the Founding-era understanding of free exercise, which protected religiously compelled conduct against even facially neutral laws unless the conduct threatened ``peace or safety'' or the rights of others. In their view, \textit{Lukumi} protects religious minorities only when majorities are clumsy or openly hostile, leaving minorities vulnerable to indifferent, generally applicable legislation.

\textbf{Progressive and Secular-Equality Perspectives:} Progressive scholars broadly celebrate \textit{Lukumi} as a vital shield for vulnerable, non-mainstream religious minorities---like Cuban-American Santeria practitioners---against majoritarian intolerance and cultural xenophobia. At the same time, progressive commentators raise two major critiques of \textit{Lukumi}'s downstream legacy. First, as the Roberts Court expanded \textit{Lukumi}'s ``general applicability'' and ``individualized exemptions'' language into a robust ``most-favored-nation'' doctrine in \textit{Fulton} and COVID-19 restriction cases, progressive critics argue that \textit{Lukumi} has been transformed from an anti-persecution shield for marginalized faiths into an aggressive deregulatory sword that allows religious claimants to demand exemptions from civil-rights, LGBTQ-nondiscrimination, and public-health laws whenever a statute contains almost any secular exception (\href{https://harvardlawreview.org/print/vol-137/rights-of-first-refusal/}{Harvard Law Review}; \href{https://scholarship.law.gwu.edu/cgi/viewcontent.cgi?filename=0&article=2806&context=faculty_publications&type=additional}{Lupu \& Tuttle on \textit{Fulton}}). Second, animal-welfare advocates note that while \textit{Lukumi} rightly condemned Hialeah's hypocrisy, its sweeping critique of underinclusiveness can make incremental animal-protection legislation harder to defend whenever legislatures regulate some harms to animals before others.

\section{Key Quotations}

\begin{quote}
Official action that targets religious conduct for distinctive treatment cannot be shielded by mere compliance with the requirement of facial neutrality. The Free Exercise Clause protects against governmental hostility which is masked as well as overt.
\begin{flushright}--- Justice Kennedy, Opinion of the Court (\href{https://supreme.justia.com/cases/federal/us/508/520/}{508 U.S. at 534})\end{flushright}
\end{quote}

\begin{quote}
The compelling interest standard that we apply once a law fails to meet the \textit{Smith} requirements is not ``water[ed] \dots{} down'' but ``really means what it says.''
\begin{flushright}--- Justice Kennedy, Opinion of the Court (\href{https://www.law.cornell.edu/supct/html/91-948.ZO.html}{508 U.S. at 546})\end{flushright}
\end{quote}

\begin{quote}
The Free Exercise Clause commits government itself to religious tolerance, and upon even slight suspicion that proposals for state intervention stem from animosity to religion or distrust of its practices, all officials must pause to remember their own high duty to the Constitution and to the rights it secures.
\begin{flushright}--- Justice Kennedy, Opinion of the Court (\href{https://www.law.cornell.edu/supct/html/91-948.ZO.html}{508 U.S. at 547})\end{flushright}
\end{quote}

\begin{quote}
The First Amendment does not refer to the purposes for which legislators enact laws, but to the effects of the laws enacted: ``Congress shall make no law \dots{} prohibiting the free exercise [of religion] \dots{}.'' This does not put us in the business of invalidating laws by reason of the evil motives of their authors.
\begin{flushright}--- Justice Scalia, Concurring in Part and Concurring in the Judgment (\href{https://supreme.justia.com/cases/federal/us/508/520/}{508 U.S. at 558})\end{flushright}
\end{quote}

\begin{quote}
``Neutral, generally applicable'' laws, drafted as they are from the perspective of the nonadherent, have the unavoidable potential of putting the believer to a choice between God and government.
\begin{flushright}--- Justice Souter, Concurring in Part and Concurring in the Judgment (\href{https://supreme.justia.com/cases/federal/us/508/520/}{508 U.S. at 577})\end{flushright}
\end{quote}

\end{document}
